%=============================================================
% Cold-Drug Benchmark Paper — Two-Column Academic Format
% Author: Buddharak Sodtana
% Advisor: Assoc. Prof. Prompong Sugunnasil
%=============================================================
\documentclass[11pt,a4paper,twocolumn]{article}

\usepackage[top=0.85in, bottom=0.85in, left=0.7in, right=0.7in]{geometry}
\setlength{\columnsep}{18pt}          % gap between the two columns

\usepackage[T1]{fontenc}
\usepackage[utf8]{inputenc}
\usepackage{setspace}
\usepackage{amsmath,amssymb}
\usepackage{booktabs}
\usepackage{array}
\usepackage{tabularx}
\usepackage{graphicx}
\usepackage{hyperref}
\usepackage{titlesec}
\usepackage{float}
\usepackage{enumitem}
\usepackage{needspace}
\usepackage{caption}
\usepackage{stfloats}                 % fixes float placement in twocolumn
\usepackage{pifont}                   % checkmark/crossmark symbols for tables
\emergencystretch=1.5em               % prevent overfull/underfull hbox in narrow columns

\hypersetup{
  colorlinks=true,
  hypertexnames=false,
  linkcolor=blue,
  urlcolor=blue,
  pdftitle={Addressing the Long-Tail in Cold-Start Drug-Drug Interaction Prediction via Graph Attention and Fingerprint Fusion},
  pdfauthor={Buddharak Sodtana}
}

\newcolumntype{Y}{>{\raggedright\arraybackslash}X}
\renewcommand{\arraystretch}{1.10}
\setstretch{1.05}
\setlist{nosep,leftmargin=*}

\titlespacing*{\section}{0pt}{10pt}{4pt}
\titlespacing*{\subsection}{0pt}{6pt}{3pt}
\titlespacing*{\subsubsection}{0pt}{4pt}{2pt}

\title{Addressing the Long-Tail in Cold-Start\\
       Drug--Drug Interaction Prediction\\
       via Graph Attention and Fingerprint Fusion}
\author{Buddharak Sodtana\\[1em]
        Advisor: Associate Professor Prompong Sugunnasil}
\date{March 2026}

\raggedbottom

\begin{document}

\maketitle

\begin{abstract}
Drug--drug interactions (DDIs) influence absorption, metabolism, efficacy, and toxicity,
making reliable prediction essential for patient safety.
However, most published evaluations allow training--test drug overlap, inflating
apparent performance and obscuring real deployment risk.
This study proposes a cold-start DDI prediction framework optimized for rare,
long-tail interaction types: all test drugs are entirely excluded from training,
reflecting the clinical scenario where newly synthesized compounds must be screened
without prior model exposure.
The framework combines a Siamese graph attention network with ECFP4 fingerprints,
logit adjustment ($\tau=0.5$), and deferred re-weighting---and is benchmarked against
an XGBoost + ECFP4 descriptor baseline under identical cold-drug folds.
Inputs are restricted to pure chemical structure (SMILES) to ensure absolute
cold-start deployment readiness, without relying on biomedical knowledge graphs
or pre-existing interaction metadata.
Across five cold-drug folds, the best neural configuration achieves mean macro
PR-AUC of 0.5519 and macro ROC-AUC of 0.9395; however, an exact paired
Wilcoxon signed-rank test finds no statistically significant difference
from the XGBoost + ECFP4 descriptor baseline (macro PR-AUC 0.5416,
all $p > 0.06$), which is also substantially better calibrated.
The defensible contribution lies within the neural family:
adding logit adjustment and deferred re-weighting to the
\texttt{GAT+ECFP} baseline raises macro PR-AUC from 0.5260 to 0.5519,
while tail-class retrieval remains difficult.
Tail-class retrieval and probability calibration remain the two
principal unresolved challenges, with staged mitigation strategies proposed for both.

\noindent\textbf{Keywords:} drug-drug interaction, cold-start, graph attention network,
ECFP, XGBoost, logit adjustment, long-tail learning, calibration
\end{abstract}

%=============================================================
\section{Introduction}
%=============================================================

\subsection{Background and Motivation}
Drug--drug interactions remain a major clinical concern because polypharmacy is common
in the treatment of chronic disease, multimorbidity, and critical illness.
Interaction-related adverse drug reactions account for approximately 6.5\% of hospital
admissions in developed countries \cite{wishart2018drugbank}, leading to reduced
therapeutic effect, and severe outcomes such as arrhythmia, hepatotoxicity, or
treatment failure.
Curated resources such as DrugBank and benchmark platforms such as Therapeutics Data
Commons (TDC) have made large-scale supervised DDI prediction increasingly practical
\cite{wishart2018drugbank, huang2021tdc}.
Despite this progress, two issues continue to limit reliable evaluation: strong
multi-class imbalance and permissive split strategies that allow test drugs to appear
during training.

\subsection{Problem Statement}
The core challenge in DDI prediction is to build a classifier that generalizes to
previously unseen compounds.
Let $d_a$ and $d_b$ denote two drugs represented by SMILES strings.
The task is to learn
\begin{equation}
  f(d_a, d_b) \rightarrow y, \qquad y \in \{1,\dots,86\},
\end{equation}
where each label represents a DrugBank interaction type.
In many standard benchmarks, a test drug can still appear in the training set in a
different pairing context, creating transductive leakage and inflating performance
estimates.
In contrast, the cold-drug setting used in this study excludes every test drug from
training entirely.
The model must therefore generalize from chemical structure alone, better reflecting
real deployment where new or infrequently used drugs must be screened without prior
model exposure.

\subsection{Scope: Structure-Only Cold-Start}
A critical framing constraint distinguishes this work from state-of-the-art models
such as EmerGNN \cite{zhang2023emergnn}, which achieves strong cold-start performance
by traversing biomedical knowledge graph paths (targets, pathways, side-effects) that
connect drugs to their biological context.
While powerful, such models cannot process entirely novel molecules that lack
established biological metadata.
This study deliberately restricts inputs to pure chemical structure
(SMILES $\to$ molecular graph / ECFP4) to guarantee absolute cold-start deployment
readiness for newly synthesized compounds that have no prior biological annotation.
This is the operationally relevant scenario for early-stage drug discovery.

\subsection{Research Objectives}
This study has four main objectives:
\begin{enumerate}
\item Build a cold-start DDI prediction framework using only chemical-structure
      inputs, with RDKit, PyTorch, and PyTorch Geometric \cite{fey2019pyg}.
\item Develop structure-aware drug representations via a Siamese graph attention
      network fused with fixed ECFP4 fingerprints,
      and compare against an XGBoost + ECFP4 descriptor baseline
      \cite{velickovic2018gat, rogers2010ecfp}.
\item Improve minority-class (long-tail) performance through logit adjustment and
      deferred re-weighting \cite{menon2021logit, cao2019ldam}.
\item Quantify the contribution of each component through a five-fold cold-drug
      ablation study, including calibration metrics alongside discrimination.
\end{enumerate}

\subsection{Contributions}
This work makes four contributions:
\begin{enumerate}
\item A documented cold-drug framework built from 190{,}996 DrugBank DDI rows with
      explicit leakage control, designed specifically around the long-tail
      interaction types that carry the highest clinical risk.
\item An apples-to-apples comparison between neural (GAT + ECFP) and non-neural
      (XGBoost + ECFP) model families under an identical strict cold-drug protocol
      on the full 86-type benchmark, providing a transparent and reproducible
      non-neural reference point not systematically available in prior work.
\item Systematic evaluation of logit adjustment and DRW under cold-drug conditions,
      showing a macro PR-AUC gain within the neural family (0.5260 $\to$ 0.5519)
      while tail-class retrieval remains an open challenge for all model families.
\item Identification of calibration degradation upon ECFP fusion as a distinct
      failure mode, with a staged mitigation roadmap (temperature scaling $\to$
      classwise scaling $\to$ conformal prediction).
\end{enumerate}

%=============================================================
\section{Literature Review}
%=============================================================

This section reviews four themes that frame the present study: the DrugBank
benchmark and its evaluation pitfalls (Section~\ref{sec:benchmark}), molecular
representations and fusion strategies (Section~\ref{sec:molrep}), the cold-drug
challenge (Section~\ref{sec:coldstart}), long-tail learning and calibration
(Section~\ref{sec:longtail}), and the positioning of this work
(Section~\ref{sec:positioning}).
Table~\ref{tab:litcomp} provides a structured comparison of representative
methods across five evaluation dimensions; the text below synthesizes cross-cutting
themes rather than enumerating individual models.

%-------------------------------------------------------------
\subsection{DDI Benchmark and Evaluation Protocols}
\label{sec:benchmark}
%-------------------------------------------------------------
The DrugBank 86-type DDI benchmark, standardized through TDC
\cite{wishart2018drugbank, huang2021tdc}, is the dominant testbed for multi-class
DDI prediction.
Under permissive (warm-start) splits, recent architectures report $>$95\% accuracy
and $>$99\% AUROC \cite{lin2024masmddi, feng2024bdnddi}, but these figures depend
critically on split design: DDI-Ben \cite{shen2025ddiben} and OpenDDI
\cite{openddi2025} show substantial degradation once test drugs diverge
structurally from training drugs, confirming that strict cold-drug protocols are
the most reliable proxy for deployment readiness
\cite{lin2023evaluation, lukashina2021coldstart}.
Lukashina et al.\ \cite{lukashina2021coldstart} formalize four cold-start levels
(T1--T4); the present study operates in the T2/T3 regime where one or both test
drugs are entirely unseen.

%-------------------------------------------------------------
\subsection{Molecular Representations and Fusion}
\label{sec:molrep}
%-------------------------------------------------------------
Two representation families dominate DDI prediction.
\textit{Graph neural networks} learn task-specific structural encodings: GAT
through multi-head attention \cite{velickovic2018gat}, GIN through expressive sum
aggregation \cite{xu2019gin}, and substructure-aware variants such as SSI-DDI
\cite{nyamabo2021ssiddi} and SSF-DDI \cite{ssfddi2024} through co-attention over
molecular fragments.
\textit{Fingerprints and gradient-boosted models} offer a complementary paradigm:
ECFP4 \cite{rogers2010ecfp} provides deterministic, portable substructure
descriptors, and XGBoost with ECFP has been shown to match or exceed GNNs on
molecular property benchmarks
\cite{boldini2023xgboost, jiang2021gnnvsdescriptor}, though its performance on
the 86-type DDI benchmark under strict cold-drug evaluation has not been
systematically reported \cite{shtar2019xgboost, tran2021xgboost}.

Fusion of these representations yields consistent gains.
Zagidullin et al.\ \cite{zagidullin2021fingerprints} show that combining
fingerprint types with low mutual similarity maximizes complementary benefit,
and graph--fingerprint fusion outperforms concatenation when cross-attention is
used \cite{mlfgnn2025}.
Knowledge-graph methods such as EmerGNN \cite{zhang2023emergnn} and KnowDDI
\cite{wang2024knowddi} achieve strong cold-start results by traversing biomedical
entity paths, but require external data infrastructure unavailable for truly novel
compounds.

%-------------------------------------------------------------
\subsection{The Cold-Drug Challenge}
\label{sec:coldstart}
%-------------------------------------------------------------
The cold-drug DDI problem requires generalizing to compounds absent from training.
Three architectural strategies have emerged.
\textit{Meta-learning} approaches such as Meta3D-DDI \cite{lv2023meta3d} apply
MAML-style bilevel optimization for scaffold-based cold-start prediction.
\textit{Substructure transfer} methods, including DSN-DDI \cite{dsnddi2023},
HDN-DDI \cite{hdnddi2025}, and the pretrained encoders of Pang et al.\ \cite{pang2023pretraining}, learn local chemical patterns that transfer to unseen
drugs.
\textit{Siamese architectures} encode each drug independently through shared
weights, making them naturally cold-start compatible
\cite{schwarz2021attentionddi, cnnsiam2023}; the present study adopts this
paradigm with graph attention.

%-------------------------------------------------------------
\subsection{Long-Tail Learning and Calibration}
\label{sec:longtail}
%-------------------------------------------------------------
DDI labels are strongly imbalanced, and the clinical stakes of missing rare types
are high: 44\% of DDI adverse event cases in the FDA FAERS database resulted in
hospitalization \cite{faers2024}, DDI-related events cause approximately 74{,}000
U.S.\ emergency visits annually \cite{ren2025rareddie}, and drugs such as
cerivastatin were withdrawn after initially rare fatal interactions went undetected
for years \cite{clinther2022ddi}.
A health equity dimension compounds the problem: pharmacogenomic variants differ
across populations \cite{pharmacogenomics2023equity}, and AI models trained on
limited demographic data may systematically miss tail-class interactions affecting
underrepresented groups \cite{habre2025aiadvancing}.

Two general-purpose techniques address class imbalance: logit adjustment shifts
decision boundaries by scaled log-priors \cite{menon2021logit}, and deferred
re-weighting (DRW) delays class-weight activation until shared representations
stabilize \cite{cao2019ldam}.
DDI-specific extensions include contrastive learning for rare events
\cite{xiong2023mrgnn}, adaptive contrastive augmentation \cite{mhafrddi2025},
metric-based meta-learning \cite{ren2025rareddie}, and the Tailed Focal Loss of
Devil-DDI \cite{devilddi2024}.
While these methods advance long-tail DDI learning, they are evaluated under
warm-start or partially overlapping splits where test drugs may appear in
training.
A systematic evaluation of logit adjustment and DRW under a strictly defined
multi-class cold-drug protocol---where all test drugs are entirely unseen and
some labels may be absent from training---across the full 86 interaction types
remains underexplored.

A related gap is the near-absence of calibration evaluation in DDI prediction.
DABA-DDI \cite{daba2025} applies Bayesian calibration, and Friesacher et al.\ \cite{friesacher2025calibration} provide a comprehensive calibration study for
structure--activity models, but no DDI study reports calibration metrics under
cold-drug conditions---a gap this work addresses.

%-------------------------------------------------------------
\subsection{Positioning of the Present Study}
\label{sec:positioning}
%-------------------------------------------------------------
Table~\ref{tab:litcomp} makes the positioning concrete.
Most existing DrugBank studies rely on permissive splits inflating metrics to
95\%+ accuracy and 99\%+ AUROC; the cold-drug protocol removes this bias.
The present study fills four specific gaps.
\textit{First}, it provides a documented cold-drug benchmark with explicit leakage
checks, positioned within the T2/T3 taxonomy \cite{lukashina2021coldstart} and
aligned with DDI-Ben's evaluation philosophy \cite{shen2025ddiben}.
\textit{Second}, while recent cold-start studies employ molecular-feature-based
models \cite{chen2023substructure, pang2023pretraining}, none include XGBoost
with ECFP4 as a direct baseline under an identical strict cold-drug protocol
on the full 86-type benchmark; this study provides that apples-to-apples
comparison \cite{boldini2023xgboost, jiang2021gnnvsdescriptor}.
\textit{Third}, it evaluates logit adjustment and DRW under cold-drug conditions,
complementing methods evaluated only under warm-start protocols.
\textit{Fourth}, it reports calibration metrics (ECE, Brier) alongside
discrimination, addressing a prerequisite for clinical DDI decision support
\cite{friesacher2025calibration}.

% Comparison table — spans both columns
\begin{table*}[tp]
\centering
\caption{Representative DDI prediction methods compared across five evaluation
         dimensions.
         \textbf{Ext.\ Bio} = requires external biomedical knowledge (targets,
         pathways, KG).
         \textbf{Cold} = evaluated under strict cold-drug protocol (all test
         drugs unseen).
         \textbf{Tail} = addresses long-tail / rare-class imbalance.
         \textbf{Cal.} = reports calibration metrics (ECE or Brier).
         \ding{51} = yes; \ding{55} = no; $\sim$ = partial.}
\label{tab:litcomp}
\resizebox{\textwidth}{!}{%
\begin{tabular}{llccccp{5.2cm}}
\toprule
Method & Architecture & Ext.\ Bio & Cold & Tail & Cal.
  & Key Limitation \\
\midrule
\multicolumn{7}{l}{\textit{Warm-start architectures}} \\
MASMDDI \cite{lin2024masmddi}
  & Soft-mask GAT       & \ding{55} & \ding{55} & \ding{55} & \ding{55}
  & 95.96\% acc.; warm split only \\
BDN-DDI \cite{feng2024bdnddi}
  & Bilinear dual-view  & \ding{55} & \ding{55} & \ding{55} & \ding{55}
  & $>$99\% AUROC; warm split only \\
SSI-DDI \cite{nyamabo2021ssiddi}
  & Substructure co-att. & \ding{55} & \ding{55} & \ding{55} & \ding{55}
  & No cold-start evaluation \\
DDIMDL \cite{deng2020ddimdl}
  & Multimodal DL        & \ding{51} & \ding{55} & \ding{55} & \ding{55}
  & Requires targets, enzymes, pathways \\[3pt]
\multicolumn{7}{l}{\textit{Cold-start architectures}} \\
EmerGNN \cite{zhang2023emergnn}
  & Flow-based GNN       & \ding{51} & \ding{51} & \ding{55} & \ding{55}
  & Requires biomedical KG paths \\
Meta3D-DDI \cite{lv2023meta3d}
  & 3D GNN + MAML        & \ding{55} & \ding{51} & \ding{55} & \ding{55}
  & Scaffold-based cold-start only \\
DSN-DDI \cite{dsnddi2023}
  & Dual-view substructure & \ding{55} & \ding{51} & \ding{55} & \ding{55}
  & No long-tail or calibration \\
SSF-DDI \cite{ssfddi2024}
  & SMILES + GNN fusion  & \ding{55} & \ding{51} & \ding{55} & \ding{55}
  & No 86-type full benchmark \\
HDN-DDI \cite{hdnddi2025}
  & Hierarchical GNN     & \ding{55} & \ding{51} & \ding{55} & \ding{55}
  & No calibration evaluation \\
SGT-DDI \cite{sgtddi2025}
  & Graph Transformer    & \ding{55} & \ding{51} & \ding{55} & \ding{55}
  & 72.81\% cold acc.; no tail or cal. \\
Pang et al.\ \cite{pang2023pretraining}
  & Pretrained encoders  & \ding{55} & \ding{51} & \ding{55} & \ding{55}
  & No XGBoost comparison \\
Chen et al.\ \cite{chen2023substructure}
  & Cross-dep.\ network  & \ding{55} & \ding{51} & \ding{55} & \ding{55}
  & Different eval.\ context \\[3pt]
\multicolumn{7}{l}{\textit{Long-tail / rare-class methods}} \\
RareDDIE \cite{ren2025rareddie}
  & Meta-learning        & \ding{55} & \ding{55} & \ding{51} & \ding{55}
  & Warm-start only \\
Devil-DDI \cite{devilddi2024}
  & Multi-view + tail loss & \ding{55} & \ding{55} & \ding{51} & \ding{55}
  & Warm-start only \\
MRCGNN \cite{xiong2023mrgnn}
  & Contrastive GNN      & \ding{55} & \ding{55} & \ding{51} & \ding{55}
  & Warm-start only \\[3pt]
\multicolumn{7}{l}{\textit{This study}} \\
\textbf{GAT+ECFP+LA+DRW}
  & Siamese GAT + ECFP4  & \ding{55} & \ding{51} & \ding{51} & \ding{51}
  & First to combine all four \\
\textbf{XGBoost+ECFP4}
  & Gradient boosting    & \ding{55} & \ding{51} & \ding{55} & \ding{51}
  & Non-neural baseline under same protocol \\
\bottomrule
\end{tabular}%
}
\end{table*}

%=============================================================
\section{Methodology}
%=============================================================

\subsection{Pipeline Overview}
The overall pipeline contains five stages: dataset loading, split generation,
graph and fingerprint construction, Siamese pair modeling, and evaluation.
The implementation uses RDKit for molecular parsing, PyTorch for optimization,
and PyTorch Geometric for graph neural network operators \cite{fey2019pyg}.
Model selection is based on validation macro PR-AUC, which is particularly
appropriate for imbalanced multi-class settings.
Labels are mapped to $\{0,\dots,85\}$ for compatibility with the 86-way
softmax classification head.

\subsection{Dataset Preparation and Cold-Drug Split Design}

\paragraph{Notation.}
Let $\mathcal{D}$ denote the set of all drugs and
$\mathcal{P} \subseteq \binom{\mathcal{D}}{2} \times \mathcal{Y}$ the set
of labelled unordered pairs, where $\mathcal{Y}=\{0,\dots,85\}$.
Each raw record $(d_i, d_j, y)$ is canonicalized into an unordered pair
$\{d_i, d_j\}$ by lexicographic ordering of the two canonical SMILES strings.

\paragraph{Conflict removal.}
Pairs that carry more than one distinct label across duplicate records are
removed entirely, eliminating label ambiguity before any model sees the data.
Of the 191{,}808 raw records, 191{,}402 unique unordered pairs are formed;
406 conflicting pairs (812 rows) are discarded, yielding a clean dataset
$\mathcal{P}^{\ast}$ of 190{,}996 interaction pairs.

\paragraph{Drug-disjoint $k$-fold partitioning.}
The drug set $\mathcal{D}$ is partitioned into $k=5$ disjoint subsets
$\mathcal{D}_0, \dots, \mathcal{D}_4$ via a degree-aware greedy balancing
procedure, such that
\begin{equation}
  \mathcal{D} = \bigsqcup_{i=0}^{4} \mathcal{D}_i, \qquad
  \mathcal{D}_i \cap \mathcal{D}_j = \emptyset \;\;\forall\; i \neq j.
\end{equation}
For each outer fold $f \in \{0,\dots,4\}$, define the test drug set
$\mathcal{D}^{\mathrm{te}}_f = \mathcal{D}_f$, the validation drug set
$\mathcal{D}^{\mathrm{val}}_f = \mathcal{D}_{(f+1)\bmod 5}$, and the
training drug set
$\mathcal{D}^{\mathrm{tr}}_f = \mathcal{D} \setminus
(\mathcal{D}^{\mathrm{te}}_f \cup \mathcal{D}^{\mathrm{val}}_f)$.

\paragraph{Pair assignment.}
Each pair $\{d_i,d_j\} \in \mathcal{P}^{\ast}$ is assigned according to
the fold membership of its constituent drugs:
\begin{equation}
\resizebox{0.88\columnwidth}{!}{$\displaystyle
  \{d_i,d_j\} \!\in\!
  \begin{cases}
    \mathcal{P}^{\mathrm{tr}}_f
      & d_i, d_j \in \mathcal{D}^{\mathrm{tr}}_f \\[2pt]
    \mathcal{P}^{\mathrm{val}}_f
      & \text{one of } d_i, d_j \in \mathcal{D}^{\mathrm{val}}_f,\;
        \text{other} \in \mathcal{D}^{\mathrm{tr}}_f \\[2pt]
    \mathcal{P}^{\mathrm{te}}_f
      & \text{one of } d_i, d_j \in \mathcal{D}^{\mathrm{te}}_f,\;
        \text{other} \in \mathcal{D}^{\mathrm{tr}}_f \\[2pt]
    \text{excluded}
      & \text{otherwise}
  \end{cases}
$}
\end{equation}
By construction, every test and validation pair contains at least one drug
whose identity is entirely absent from the training drug set, enforcing the
cold-drug constraint
$\mathcal{D}^{\mathrm{tr}}_f \cap \mathcal{D}^{\mathrm{te}}_f = \emptyset$.
This corresponds to a T2-like scenario (one novel drug paired with one
training-set drug), which provides a controlled evaluation of
generalization to unseen compounds.
Pairs where both drugs belong to holdout folds are explicitly excluded:
retaining them would conflate T3-level difficulty (two novel drugs) with
T2 and risk transductive leakage between the validation and test partitions.
Mean split statistics across folds and the strongest fold are shown in
Table~\ref{tab:split_stats}.

\begin{table}[H]
\centering
\caption{Cold-drug split statistics.}
\label{tab:split_stats}
\resizebox{\columnwidth}{!}{%
\begin{tabular}{lrr}
\toprule
Statistic & Mean (5 Folds) & Fold $f_4$ \\
\midrule
Training pairs $|\mathcal{P}^{\mathrm{tr}}|$
                         & 68{,}619.4 & 68{,}562 \\
Validation pairs $|\mathcal{P}^{\mathrm{val}}|$
                         & 45{,}978.2 & 46{,}038 \\
Test pairs $|\mathcal{P}^{\mathrm{te}}|$
                         & 45{,}978.2 & 46{,}033 \\
Excluded pairs           & 30{,}420.2 & 30{,}363 \\
Test labels present      & 81.8       & 85       \\
Test labels absent       & 4.2        & 3        \\
Unique drugs assigned    & 341.2      & 342      \\
\bottomrule
\end{tabular}%
}
\end{table}

Partition sizes do not sum to $|\mathcal{P}^{\ast}| = 190{,}996$ because
pairs connecting two holdout drugs are excluded for each outer fold.
The drug-disjoint partitioning does not guarantee full label coverage:
depending on the fold, 3--6 of the 86 interaction types are absent from
the training partition.
This label-coverage gap is reported explicitly rather than hidden.

\subsection{Molecular Graph and Fingerprint Features}
Each drug is converted from SMILES into a molecular graph using RDKit
(invalid or zero-atom molecules are skipped).
Bonds are stored as directed edges in both directions; self-loops are added
internally by each GAT layer.
Table~\ref{tab:features} summarizes the node and edge feature sets.

\begin{table}[H]
\centering
\caption{Molecular graph feature summary.}
\label{tab:features}
\resizebox{\columnwidth}{!}{%
\begin{tabular}{llc}
\toprule
Scope & Feature & Values \\
\midrule
Node (7\,dim) & Atomic number, degree, formal charge,
                 total H count & integer \\
              & Aromaticity, ring membership & 0/1 \\
              & Hybridization & 0--7 ordinal \\
Edge (5\,dim) & Bond type (single/double/triple/aromatic) & 1--4 \\
              & Conjugation, aromaticity, ring membership & 0/1 \\
              & Bond stereo & 0--5 ordinal \\
\bottomrule
\end{tabular}%
}
\end{table}

ECFP4 fingerprints (radius~2, 2048 bits) \cite{rogers2010ecfp} are computed
with RDKit Morgan fingerprints (\texttt{useChirality=False}) and projected via
$\texttt{Linear}(2048,128) \!\to\! \texttt{LayerNorm} \!\to\!
\texttt{ReLU} \!\to\! \texttt{Dropout}(0.2)$
before fusion with the graph embedding.

\subsection{Siamese Pair Model}
The model uses a Siamese architecture in which both drugs are encoded by the
same graph attention backbone.
The encoder contains three GAT layers with edge-aware attention and mean graph
pooling.
The first two layers use 64 hidden channels per head, 4 heads, concatenated
multi-head outputs, ELU activation, and dropout~0.2.
The final layer uses 128 output channels with 4 heads and \texttt{concat=False},
producing a 128-dimensional pooled graph embedding~$\mathbf{g}$.
Let $\mathbf{g}_a$ and $\mathbf{g}_b$ denote the graph embeddings for drugs~$a$
and~$b$, and let $\mathbf{p}_a$ and $\mathbf{p}_b$ denote the projected
fingerprint vectors.
Because the strongest baseline enables the ECFP pathway, each drug embedding is
computed by
\begin{equation}
  \mathbf{h}_a = \phi([\mathbf{g}_a \| \mathbf{p}_a]), \qquad
  \mathbf{h}_b = \phi([\mathbf{g}_b \| \mathbf{p}_b]),
\end{equation}
where $\phi$ is a fusion block
$\texttt{Linear}(256,128) \rightarrow \texttt{ReLU} \rightarrow
\texttt{Dropout}(0.2)$.

The pairwise representation is permutation-invariant in implementation.
Define element-wise $\mathbf{h}_{\min} = \min(\mathbf{h}_a,\mathbf{h}_b)$ and
$\mathbf{h}_{\max} = \max(\mathbf{h}_a,\mathbf{h}_b)$.
Note that $\mathbf{h}_{\max} - \mathbf{h}_{\min} = |\mathbf{h}_a - \mathbf{h}_b|$
element-wise.
The pair feature is therefore
\begin{equation}
  \mathbf{z}_{ab} =
  [\mathbf{h}_{\min} \| \mathbf{h}_{\max} \|
  |\mathbf{h}_a - \mathbf{h}_b| \|
  \mathbf{h}_a \odot \mathbf{h}_b],
\end{equation}
where $|\cdot|$ is the element-wise absolute difference.
With $\mathbf{h}_a, \mathbf{h}_b \in \mathbb{R}^{128}$, the resulting pair
feature dimension is~512.
A classifier maps this representation through
$\texttt{Linear}(512,256) \rightarrow \texttt{ReLU} \rightarrow
\texttt{Dropout}(0.2) \rightarrow \texttt{Linear}(256,86)$
to produce the final logits.

\subsection{XGBoost + ECFP4 Descriptor Baseline}
\label{sec:xgboost}
To provide a strong non-neural reference under the same cold-drug folds, the
benchmark includes an XGBoost multiclass classifier trained solely on ECFP4
pair descriptors.
This baseline tests how much performance is attainable from fixed substructure
fingerprints without learned graph message passing, and establishes a
structural lower bound for any fair architectural comparison.

Let $\mathbf{f}(d)\!\in\!\{0,1\}^{2048}$ denote the ECFP4 fingerprint.
Pairs are lexicographically sorted by canonical SMILES, yielding a
permutation-consistent pair vector
$\mathbf{x}_{ab} = [\mathbf{f}(d_{(1)}) \| \mathbf{f}(d_{(2)})]
\in \{0,1\}^{4096}$.

The XGBoost classifier uses the multiclass soft-probability objective
(\texttt{multi:softprob}) with histogram-based tree growth (\texttt{tree\_method=hist}).
Class weighting is intentionally omitted to establish a fair
structural-feature-only baseline without any long-tail intervention.
Because some folds contain labels absent from training, XGBoost trains only on
the labels observed in the training partition; validation rows whose labels are
absent from training are dropped from the early-stopping evaluation set, while
test predictions are expanded to the full 86-class space by assigning zero
probability to train-unseen classes.
The full hyperparameter configuration is listed in
Table~\ref{tab:xgboost_hyperparams}.

\subsection{Long-Tail Objective}
The model uses logit adjustment based on empirical train-split priors computed
inside each fold with $\epsilon = 10^{-12}$.
Let $\mathbf{s}$ denote the raw classifier logits and $\boldsymbol{\pi}$ the
class-prior vector.
The adjusted logits are
\begin{equation}
  \tilde{\mathbf{s}} = \mathbf{s} + \tau \log \boldsymbol{\pi},
\end{equation}
where $\tau = 0.5$ \cite{menon2021logit}.
Cross-entropy is then applied to the adjusted logits, optionally with class
weights:
\begin{equation}
  \mathcal{L} = \mathrm{CE}(\tilde{\mathbf{s}}, y; \mathbf{w}).
\end{equation}

During epochs~1--14, the model trains without class weighting.
From epoch~15 onward, deferred re-weighting is activated and class weights are
defined by
\begin{equation}
  w_c = \mathrm{clip}\!\left(\frac{1}{\sqrt{n_c + \epsilon}},\;
  0.25,\; 4.0\right),
\end{equation}
where $n_c$ is the number of training examples in class~$c$ and
$\epsilon = 10^{-12}$.
In implementation, seen-class weights are normalized before clipping and are
applied only after the DRW switch.
At the same point, the learning rate is reduced from $10^{-3}$ to
$2 \times 10^{-4}$.
This schedule reflects the DRW intuition that shared feature learning should
stabilize before minority-focused re-weighting is introduced \cite{cao2019ldam}.

\subsection{Experimental Configuration}
The final neural configuration was chosen by a bounded development-fold search,
not by tuning on the full five-fold benchmark.
No outer-fold test set was used at any stage of model selection.

\paragraph{Hyperparameter selection.}
A predefined grid ($\tau \!\in\! \{0.4, 0.5, 0.75\}$,
LR $\!\in\! \{3\!\times\!10^{-4}\!,\;5\!\times\!10^{-4}\!,\;10^{-3}\!,\;2\!\times\!10^{-3}\}$,
DRW ratio $\!\in\! \{0.6, 0.7, 0.8\}$) was screened on a single development
fold (fold~0, training cap 100{,}000 rows) ranked by
tail PR-AUC $\succ$ macro PR-AUC $\succ$ macro-F1 $\succ$ lower loss.
The top-ranked configuration was $\tau\!=\!0.5$, LR $10^{-3}$, DRW ratio~0.7
(start epoch~15).
Architectural settings (3 GAT layers, hidden 64, output 128, 4 heads,
dropout 0.2) were held fixed and are not claimed to be globally optimal;
the protocol is designed for a reproducible cold-drug reference under a
fixed compute budget.
The XGBoost configuration was likewise fixed without outer-fold tuning.
Neural and XGBoost hyperparameters are summarized in
Tables~\ref{tab:hyperparams} and~\ref{tab:xgboost_hyperparams}.

\begin{table}[H]
\centering
\caption{Core model and training hyperparameters.}
\label{tab:hyperparams}
\resizebox{\columnwidth}{!}{%
\begin{tabular}{ll}
\toprule
Component & Configuration \\
\midrule
Graph encoder            & GAT \\
Node / edge feature dim  & 7 / 5 \\
Number of GAT layers     & 3 \\
Attention heads          & 4 \\
Hidden dimension         & 64 \\
Graph embedding dim      & 128 \\
Pooling                  & global mean \\
Fingerprint              & ECFP4, r=2, 2048 bits \\
Fingerprint proj.\ dim   & 128 \\
Pair feature dim         & 512 \\
Classifier hidden dim    & 256 \\
Dropout                  & 0.2 \\
Batch size / epochs      & 64 / 20 (max) \\
Optimizer                & Adam \\
Learning rate            & $10^{-3}$, then $2\!\times\!10^{-4}$ after DRW \\
Weight decay             & $10^{-5}$ \\
Early stopping           & patience 5 on val.\ macro PR-AUC \\
Logit adjustment         & $\tau = 0.5$ \\
DRW ratio                & 0.7 of training horizon \\
Selection metric         & validation macro PR-AUC \\
\bottomrule
\end{tabular}%
}
\end{table}

\begin{table}[htbp]
\centering
\caption{XGBoost baseline hyperparameters.}
\label{tab:xgboost_hyperparams}
\resizebox{\columnwidth}{!}{%
\begin{tabular}{ll}
\toprule
Component & Configuration \\
\midrule
Feature type           & sorted concat.\ ECFP4 \\
Feature dimension      & 4{,}096 \\
ECFP radius / bits     & 2 / 2{,}048 \\
Objective              & \texttt{multi:softprob} \\
\texttt{n\_estimators} & 600 \\
Learning rate          & 0.05 \\
Max depth              & 8 \\
Subsample              & 0.8 \\
\texttt{colsample\_bytree} & 0.8 \\
\texttt{min\_child\_weight} & 1.0 \\
$\lambda$              & 1.0 \\
Tree method            & \texttt{hist} \\
Early stopping rounds  & 30 \\
Class weights          & none \\
\bottomrule
\end{tabular}%
}
\end{table}

\subsection{Evaluation Metrics}
All random seeds are fixed for reproducibility; the best checkpoint is
selected by validation macro PR-AUC.
Evaluation includes accuracy, macro-F1, Cohen's kappa, macro ROC-AUC,
macro PR-AUC, ECE (15 equal-width bins on max-softmax confidence), and
multiclass Brier score.
Macro ROC-AUC and macro PR-AUC are one-vs-rest averages over classes with
both positives and negatives.
Classes with zero training support in a given fold (3--6 of 86) are excluded
from that fold's macro-averaged evaluations.
Tail macro PR-AUC is computed over the bottom
$\lceil 0.2 \times 86 \rceil = 18$ classes ranked by training support,
including zero-count classes.

%=============================================================
\section{Results and Discussion}
%=============================================================

\subsection{Five-Fold Cold-Drug Performance}
Table~\ref{tab:main_results} reports mean $\pm$ standard deviation across five
cold-drug folds for all models, including the XGBoost + ECFP4 descriptor baseline.
The results tell a two-step story about where neural graph representations add value.

\textit{Step~1 --- Architectural advantage.}
Comparing unweighted baselines, \texttt{GAT+ECFP} (no class weighting) reaches
macro ROC-AUC 0.9286 and kappa 0.5524, exceeding the unweighted \texttt{XGBoost+ECFP4}
(ROC-AUC 0.9230, kappa 0.5310) under the same cold-drug protocol.
Both models receive ECFP4 fingerprints, but the neural model additionally
learns task-specific atomic neighborhood representations through the
graph-attention layers.
This apples-to-apples comparison therefore demonstrates that the Siamese
graph-attention mechanism captures relational structural features that
fixed concatenated fingerprints alone cannot.

\textit{Step~2 --- Long-tail mitigation.}
The best-performing neural configuration (\texttt{GAT+ECFP+LA+DRW}) achieves the
highest macro PR-AUC (0.5519), macro-F1 (0.5032), accuracy (0.6351),
macro ROC-AUC (0.9395), and kappa (0.5579).
Tail macro PR-AUC remains low at 0.2156, indicating that rare-class retrieval
is still the main unresolved difficulty.
The neural model is slightly ahead of XGBoost on mean accuracy, macro-F1,
macro PR-AUC, macro ROC-AUC, and kappa, while XGBoost is slightly higher on
mean tail macro PR-AUC (0.2174 vs.\ 0.2156) and substantially better calibrated
(ECE 0.0479 vs.\ 0.1966), a 4$\times$ advantage discussed further in
Section~\ref{sec:calib_mech}.

% Wide table: spans both columns
\begin{table*}[htbp]
\centering
\caption{Five-fold cold-drug performance. Mean $\pm$ std. Bold = best.
         Lower is better for ECE and Brier. $\dagger$~no class weighting.}
\label{tab:main_results}
\resizebox{\textwidth}{!}{%
\begin{tabular}{lcccccccc}
\toprule
Model & Accuracy & Macro-F1 & Kappa & Macro ROC-AUC & Macro PR-AUC
      & Tail PR-AUC & ECE $\downarrow$ & Brier $\downarrow$ \\
\midrule
GAT only$^\dagger$
  & 0.5927{\scriptsize$\pm$.003} & 0.4212{\scriptsize$\pm$.012}
  & 0.5030{\scriptsize$\pm$.003} & 0.9173{\scriptsize$\pm$.004}
  & 0.4616{\scriptsize$\pm$.020} & 0.1648{\scriptsize$\pm$.040}
  & 0.1410{\scriptsize$\pm$.022} & 0.5827{\scriptsize$\pm$.011} \\
GIN only$^\dagger$
  & 0.5757{\scriptsize$\pm$.008} & 0.4066{\scriptsize$\pm$.028}
  & 0.4797{\scriptsize$\pm$.009} & 0.9103{\scriptsize$\pm$.008}
  & 0.4463{\scriptsize$\pm$.028} & 0.1609{\scriptsize$\pm$.029}
  & \textbf{0.1160{\scriptsize$\pm$.009}} & 0.5945{\scriptsize$\pm$.012} \\
GAT+ECFP$^\dagger$
  & 0.6286{\scriptsize$\pm$.010} & 0.4825{\scriptsize$\pm$.030}
  & 0.5524{\scriptsize$\pm$.011} & 0.9286{\scriptsize$\pm$.009}
  & 0.5260{\scriptsize$\pm$.029} & 0.2142{\scriptsize$\pm$.063}
  & 0.1860{\scriptsize$\pm$.016} & 0.5642{\scriptsize$\pm$.013} \\
XGBoost+ECFP4$^\dagger$
  & 0.6269{\scriptsize$\pm$.007} & 0.4941{\scriptsize$\pm$.032}
  & 0.5310{\scriptsize$\pm$.007} & 0.9230{\scriptsize$\pm$.013}
  & 0.5416{\scriptsize$\pm$.039} & \textbf{0.2174{\scriptsize$\pm$.098}}
  & \textbf{0.0479{\scriptsize$\pm$.009}} & \textbf{0.5243{\scriptsize$\pm$.009}} \\
GAT+ECFP+LA
  & 0.6276{\scriptsize$\pm$.015} & 0.4915{\scriptsize$\pm$.026}
  & 0.5483{\scriptsize$\pm$.019} & 0.9352{\scriptsize$\pm$.012}
  & 0.5316{\scriptsize$\pm$.031} & 0.2163{\scriptsize$\pm$.067}
  & 0.1963{\scriptsize$\pm$.016} & 0.5704{\scriptsize$\pm$.022} \\
GAT+ECFP+LA+DRW
  & \textbf{0.6351{\scriptsize$\pm$.012}} & \textbf{0.5032{\scriptsize$\pm$.034}}
  & \textbf{0.5579{\scriptsize$\pm$.015}} & \textbf{0.9395{\scriptsize$\pm$.010}}
  & \textbf{0.5519{\scriptsize$\pm$.039}} & 0.2156{\scriptsize$\pm$.081}
  & 0.1966{\scriptsize$\pm$.019} & 0.5643{\scriptsize$\pm$.016} \\
\bottomrule
\end{tabular}%
}
\end{table*}

Table~\ref{tab:fold_results} shows fold-wise results for the best configuration.
Performance is stable: macro ROC-AUC ranges from 0.9215 to 0.9519, and accuracy from
0.6227 to 0.6563.
Fold $f4$ achieves notably higher tail PR-AUC (0.3740 vs.\ mean 0.2156), explained
by stronger class coverage (85 of 86 classes vs.\ average 81.8).

% Wide table: spans both columns
\begin{table*}[tp]
\centering
\caption{Fold-wise test performance of GAT+ECFP+LA+DRW.}
\label{tab:fold_results}
\resizebox{0.88\textwidth}{!}{%
\begin{tabular}{lrrrrrr}
\toprule
Fold & Accuracy & Macro-F1 & Macro PR-AUC & Tail PR-AUC & Macro ROC-AUC & Kappa \\
\midrule
$f0$ & 0.6325 & 0.4813 & 0.5643 & 0.1888 & 0.9464 & 0.5561 \\
$f1$ & 0.6251 & 0.4782 & 0.5101 & 0.1912 & 0.9368 & 0.5446 \\
$f2$ & 0.6388 & 0.5148 & 0.5260 & 0.1790 & 0.9409 & 0.5582 \\
$f3$ & 0.6227 & 0.4770 & 0.5371 & 0.1448 & 0.9215 & 0.5454 \\
$f4$ & 0.6563 & 0.5644 & 0.6220 & 0.3740 & 0.9519 & 0.5851 \\
\midrule
Mean & 0.6351 & 0.5032 & 0.5519 & 0.2156 & 0.9395 & 0.5579 \\
Std. & 0.0120 & 0.0337 & 0.0392 & 0.0809 & 0.0103 & 0.0147 \\
\bottomrule
\end{tabular}%
}
\end{table*}

\subsection{Descriptor Baseline Comparison}
\label{sec:xgb_comparison}
The XGBoost baseline is not a trivial comparator.
Under the strict cold-drug protocol, it outperforms the unweighted graph-only
neural baselines on macro PR-AUC and remains close to the final
\texttt{GAT+ECFP+LA+DRW} configuration across all discrimination metrics.
The mean differences are modest: +0.82~pp in accuracy, +0.91~pp in macro-F1,
+1.03~pp in macro PR-AUC, and +1.65~pp in macro ROC-AUC---while XGBoost is
slightly higher on tail PR-AUC (0.2174 vs.\ 0.2156) and
four times better calibrated (ECE 0.0479 vs.\ 0.1966).
This confirms that a substantial fraction of the cold-drug signal is already
captured by sparse ECFP4 descriptors plus a high-capacity tree ensemble.

\paragraph{Statistical significance.}
To test whether these fold-level differences are robust, an exact paired
Wilcoxon signed-rank test was computed across the five fold pairs
$f_0$--$f_4$.
Table~\ref{tab:wilcoxon_xgb} reports the two-sided exact $p$-values.
None of the comparisons reaches the conventional $p < 0.05$ threshold.
The tail-class result is the most telling: XGBoost is slightly higher on
five-fold mean tail macro PR-AUC (0.2174 vs.\ 0.2156), and the paired
Wilcoxon test gives $p = 1.0000$.
The current benchmark therefore does \emph{not} establish a statistically
significant tail-class advantage for the neural model over XGBoost.

\begin{table}[H]
\centering
\caption{Exact paired Wilcoxon signed-rank test:
         \texttt{GAT+ECFP+LA+DRW} vs.\ \texttt{XGBoost+ECFP4}
         across five cold-drug folds.
         Lower is better for ECE and Brier.}
\label{tab:wilcoxon_xgb}
\resizebox{\columnwidth}{!}{%
\begin{tabular}{lccc}
\toprule
Metric & Neural Mean & XGB Mean & Exact $p$ \\
\midrule
Accuracy         & 0.6351 & 0.6269 & 0.4375 \\
Macro-F1         & 0.5032 & 0.4941 & 0.6250 \\
Macro PR-AUC     & 0.5519 & 0.5416 & 0.3125 \\
Tail PR-AUC      & 0.2156 & 0.2174 & 1.0000 \\
Macro ROC-AUC    & 0.9395 & 0.9230 & 0.0625 \\
Kappa            & 0.5579 & 0.5310 & 0.0625 \\
ECE $\downarrow$ & 0.1966 & 0.0479 & 0.0625 \\
Brier $\downarrow$ & 0.5643 & 0.5243 & 0.0625 \\
\bottomrule
\end{tabular}%
}
\end{table}

Three limitations of this analysis should be stated.
First, with only five paired folds the exact Wilcoxon test is underpowered:
even if one model won all five folds, the smallest attainable exact two-sided
$p$-value is 0.0625.
Second, the folds are repeated holdout configurations from the same benchmark,
not independent datasets.
Third, fold~$f_0$ was also used during the development-fold hyperparameter
search, so the resulting $p$-values should be interpreted as exploratory
rather than confirmatory.

\paragraph{Where the neural model's contribution lies.}
Given the non-significant cross-family comparison, the more defensible
architectural claim is \emph{within} the neural family.
Within the neural family, logit adjustment and deferred re-weighting improve
overall retrieval more clearly than they rescue the hardest tail classes.
Relative to \texttt{GAT+ECFP}, macro PR-AUC rises from 0.5260 to 0.5519,
while tail PR-AUC changes only marginally from 0.2142 to 0.2156.
This effect is absent from the XGBoost baseline, which has
no long-tail intervention.
The XGBoost comparison therefore serves a different purpose: it establishes a
strong, well-calibrated reference point that future cold-drug models must
exceed on \emph{both} ranking and calibration quality.

\paragraph{Clinical implications of the calibration gap.}
The four-fold ECE difference carries direct clinical consequences.
An overconfident model can assign high probability to an incorrect ``safe''
label, leading clinicians to overlook a genuinely hazardous interaction;
conversely, it can produce confidently wrong positive predictions that trigger
unnecessary treatment modifications.
In high-stakes DDI screening, the reliability of predicted confidence is at
least as important as the ranking of candidate interactions, because clinical
decision support systems must communicate uncertainty faithfully to the
prescribing physician.
The more trustworthy confidence estimates from the XGBoost baseline may
therefore be preferable for deployment despite comparable ranking performance,
underscoring that calibration is a non-negotiable prerequisite for trustworthy
AI-assisted medication safety.

\subsection{Training Dynamics and DRW Behavior}
The best validation checkpoint occurs at epoch~18 (accuracy 0.6390, macro-F1 0.5246,
macro PR-AUC 0.5678, macro ROC-AUC 0.9417).
Performance improves rapidly in the first five epochs, stabilizes during the
uniform-weighting phase (epochs~1--14), and continues to improve after DRW activates
at epoch~15 without any drop at the transition point, supporting the DRW rationale
that shared representations should mature before minority-class emphasis is introduced
\cite{cao2019ldam}.

\subsection{Ablation: Effect of ECFP Fusion}
Table~\ref{tab:ecfp_ablation} isolates the contribution of fingerprint fusion.

% Wide table: spans both columns
\begin{table*}[tp]
\centering
\caption{Ablation: ECFP fusion under identical logit adjustment and DRW.
         Mean across five folds.}
\label{tab:ecfp_ablation}
\resizebox{0.90\textwidth}{!}{%
\begin{tabular}{lrrrrrrr}
\toprule
Model & Accuracy & Macro-F1 & Macro PR-AUC & Tail PR-AUC
      & Macro ROC-AUC & Kappa & ECE \\
\midrule
GAT + LA + DRW
  & 0.5967 & 0.4569 & 0.4908 & 0.2142 & 0.9259 & 0.5085 & 0.1394 \\
GAT + ECFP + LA + DRW
  & 0.6351 & 0.5032 & 0.5519 & 0.2156 & 0.9395 & 0.5579 & 0.1966 \\
\midrule
Absolute gain
  & +0.0383 & +0.0463 & +0.0610 & +0.0014 & +0.0137 & +0.0493 & +0.0572 \\
\bottomrule
\end{tabular}%
}
\end{table*}

ECFP fusion yields substantial discrimination gains within the neural family
(accuracy +3.8~pp, macro PR-AUC +12--14\%), but calibration degrades
significantly (ECE: 0.1394 $\to$ 0.1966, a relative increase of 41\%).
This result is consistent with the cross-family comparison: the XGBoost baseline
uses only ECFP4 descriptors and achieves both strong ranking (macro PR-AUC 0.5416)
\emph{and} excellent calibration (ECE 0.0479), showing that ECFP information is
powerful but the neural model's way of fusing it introduces overconfidence.
The discrimination--calibration trade-off upon feature fusion is a critical yet
rarely quantified failure mode in empirical DDI prediction.
High-performing recent architectures such as MASMDDI \cite{lin2024masmddi} and
BDN-DDI \cite{feng2024bdnddi} report accuracy and AUROC under permissive splits
but do not evaluate calibration; even comprehensive reviews of DDI methods
\cite{daba2025} typically omit ECE or Brier metrics.
While Friesacher et al.\ \cite{friesacher2025calibration} provide a thorough
calibration study for structure--activity models in general, the specific
observation that ECFP fusion degrades calibration in a multi-class DDI setting
under cold-start conditions has not been documented previously.
This gap underscores the need for post-hoc calibration as a mandatory
deployment step whenever fingerprint fusion is employed.

\subsection{Calibration Degradation: Mechanisms and Evidence}
\label{sec:calib_mech}
The ECE increase upon ECFP fusion (0.1394 $\to$ 0.1966) is a predictable consequence
of three interacting mechanisms.
First, cross-entropy training with hard labels drives classifiers toward large logit
margins; ECFP fusion raises effective model capacity, enabling greater perceived
separability---and thus higher confidence---faster than genuine correctness gains,
a global overconfidence pattern documented across model families when capacity is
scaled without re-calibration \cite{friesacher2025calibration}.
Second, both graph embeddings and ECFP4 fingerprints are derived from the same 2D
molecular topology; fusing correlated representations allows the model to count the
same structural evidence twice, amplifying confidence without a proportional accuracy
increase (fingerprint types with low Centered Kernel Alignment similarity yield more
complementary fusion gains \cite{zagidullin2021fingerprints}).
Third, calibration is distribution-dependent: a model calibrated on i.i.d.\ validation
data degrades under chemical novelty of cold-drug test sets, where substructure
patterns learned at training time may not transfer reliably.
Standard accuracy and ROC-AUC do not surface these failure modes; ECE and Brier
score are essential complements for deployment-relevant evaluation.
From a patient-safety perspective, any model deployed for DDI screening must
communicate uncertainty faithfully: a miscalibrated prediction that appears
confident can suppress clinical vigilance for dangerous interactions or
trigger unnecessary interventions for benign ones.
The four-fold ECE gap between the neural model and XGBoost therefore
represents not merely a statistical shortcoming but a concrete barrier to
clinical adoption that must be resolved---through post-hoc recalibration or
architectural redesign---before the neural model's ranking advantage
translates into real-world patient benefit.

\subsection{Tail-Class Performance}
Tail macro PR-AUC averages only 0.2156 for the best neural system, substantially
below the overall macro PR-AUC of 0.5519; XGBoost reaches 0.2174, statistically
indistinguishable within the reported variance.
Both model families face the same structural difficulty: the cold-drug protocol
removes transductive signal, compounding the scarcity of training examples for
rare classes, and some tail labels are absent from training entirely in certain
folds (Section~3.2), limiting correction by logit adjustment.
Cold-drug DDI prediction at the macro level (ROC-AUC $>$0.92) is feasible, but
tail-class retrieval and calibration remain the two open bottlenecks, for neural
and descriptor-based models alike.

%=============================================================
\section{Conclusion and Future Work}
%=============================================================

This study proposed a cold-start DDI prediction framework targeting rare,
long-tail interaction types.
Inputs are restricted to pure chemical structure (SMILES $\to$ graph / ECFP4),
guaranteeing deployment readiness for novel compounds without biological metadata.
Across five cold-drug folds, the best neural configuration (\texttt{GAT+ECFP+LA+DRW})
achieves macro PR-AUC 0.5519, macro-F1 0.5032, accuracy 0.6351, and macro
ROC-AUC 0.9395.
An important negative result is that graph neural networks do not dominate
once the evaluation is strict: an exact paired Wilcoxon signed-rank test
across the five folds finds no statistically significant difference between
the neural model and the XGBoost + ECFP4 baseline on any discrimination
metric (all $p > 0.06$), while XGBoost is substantially better calibrated
(ECE 0.0479 vs.\ 0.1966).
This makes XGBoost a mandatory comparator for any future cold-drug DDI study.
The defensible contribution of this work lies within the neural family:
adding logit adjustment and deferred re-weighting to the \texttt{GAT+ECFP}
baseline raises macro PR-AUC from 0.5260 to 0.5519, while tail-class
retrieval remains difficult (tail PR-AUC 0.2156)---underscoring that
rare interaction types under chemical novelty remain the principal
open challenge.

Three challenges remain open.
\textit{First}, tail-class retrieval is weak (mean tail macro PR-AUC 0.2156 for the
neural system and 0.2174 for XGBoost), and future work should pursue three
concrete directions.
\textit{(a)~Metric learning for cold-drug tail classes.}
Contrastive or triplet loss objectives, analogous to those used by Karim et al.\
\cite{karim2019kgddi} for knowledge-graph-based DDI embeddings, could be adapted
to learn drug representations in which rare interaction types occupy more
discriminative regions of the embedding space; the key challenge under the
cold-drug constraint is that the anchor--positive--negative sampling strategy
must generalize from training-set drugs to structurally novel test compounds.
\textit{(b)~Tail-aware augmentation and loss design.}
The Tailed Focal Loss introduced by Devil-DDI \cite{devilddi2024} dynamically
adjusts per-class focus based on both frequency and prediction difficulty;
combining such a loss with the DRW schedule used in this study could yield
complementary tail-class gains, though the interaction between deferred
re-weighting and adaptive focal terms requires careful ablation under
cold-drug splits.
\textit{(c)~Hierarchical label modeling.}
The 86 DrugBank interaction types can be grouped by pharmacological mechanism
(e.g., enzyme inhibition, transporter competition, receptor antagonism);
a hierarchical classifier that first predicts the mechanism group and then
refines to the specific type could transfer head-class knowledge to
structurally related tail classes---an approach explored in general long-tail
recognition \cite{cao2019ldam} but not yet systematically applied to
cold-start DDI prediction.

\textit{Second}, calibration worsens when ECFP fusion is added (ECE 0.1394
$\to$ 0.1966, a 41\% relative increase) due to increased capacity, correlated
evidence double-counting, and cold-drug distributional shift (Section~4.4).
Mitigation should proceed in stages: post-hoc temperature scaling on fused logits
is the required first step, as it preserves argmax predictions while reducing
global overconfidence using a single parameter; if classwise analysis reveals
heterogeneous miscalibration across the 86 DDI types, vector or classwise
temperature scaling should follow; and for deployment under chemical novelty,
where post-hoc scaling may degrade, conformal prediction provides
distribution-free coverage guarantees as a fallback \cite{friesacher2025calibration}.

\textit{Third}, future work may incorporate richer modalities such as target
proteins, pathway information, or biomedical knowledge graphs to improve prediction
for rare or novel interactions.
Interpretability methods, including saliency maps and attention visualizations,
should also be evaluated on the same cold-drug baseline so that explanation claims
rest on a reliable predictive foundation.

%=============================================================
\needspace{5\baselineskip}
%=============================================================

\section*{References}
\addcontentsline{toc}{section}{References}

\begin{thebibliography}{99}
\small
\sloppy
\setlength{\itemsep}{0pt}
\setlength{\parsep}{0pt}
\setlength{\parskip}{0pt}

\bibitem{wishart2018drugbank}
D.S.\ Wishart, Y.D.\ Feunang, A.C.\ Guo, et al.,
DrugBank 5.0: a major update to the DrugBank database for 2018,
\emph{Nucleic Acids Res.} 46 (D1) (2018) D1074--D1082.

\bibitem{huang2021tdc}
K.\ Huang, T.\ Fu, W.\ Gao, et al.,
Therapeutics Data Commons: ML datasets and tasks for therapeutics,
in: \emph{NeurIPS Datasets and Benchmarks}, 2021.

\bibitem{velickovic2018gat}
P.\ Veli\v{c}kovi\'{c}, G.\ Cucurull, A.\ Casanova, et al.,
Graph attention networks,
in: \emph{ICLR}, 2018.

\bibitem{rogers2010ecfp}
D.\ Rogers, M.\ Hahn,
Extended-connectivity fingerprints,
\emph{J.\ Chem.\ Inf.\ Model.} 50 (5) (2010) 742--754.

\bibitem{xu2019gin}
K.\ Xu, W.\ Hu, J.\ Leskovec, et al.,
How powerful are graph neural networks?,
in: \emph{ICLR}, 2019.

\bibitem{fey2019pyg}
M.\ Fey, J.E.\ Lenssen,
Fast graph representation learning with PyTorch Geometric,
in: \emph{ICLR Workshop on Representation Learning on Graphs and Manifolds}, 2019.

\bibitem{menon2021logit}
A.K.\ Menon, S.\ Jayasumana, A.S.\ Rawat, et al.,
Long-tail learning via logit adjustment,
in: \emph{ICLR}, 2021.

\bibitem{cao2019ldam}
K.\ Cao, C.\ Wei, A.\ Gaidon, et al.,
Learning imbalanced datasets with label-distribution-aware margin loss,
in: \emph{NeurIPS}, vol.\ 32, 2019.

\bibitem{xiong2023mrgnn}
Z.\ Xiong, F.\ Liu, S.\ Wang, et al.,
Multi-relational contrastive learning GNN for DDI event prediction,
in: \emph{Proc.\ AAAI}, 2023.

\bibitem{lin2023evaluation}
X.\ Lin, Z.\ Quan, Z.J.\ Wang, et al.,
Comprehensive evaluation of deep and graph learning on drug--drug interaction prediction,
\emph{Brief.\ Bioinform.} 24 (4) (2023) bbad235.

\bibitem{lin2024masmddi}
S.\ Lin, Y.\ Chen, Z.\ Li, et al.,
MASMDDI: multi-layer adaptive soft-mask graph neural network for DDI prediction,
\emph{Front.\ Pharmacol.} 15 (2024) 1369403.

\bibitem{feng2024bdnddi}
X.\ Feng, Y.\ Zhao, Z.\ Wang, et al.,
BDN-DDI: a bilinear dual-view representation learning framework for DDI prediction,
\emph{Comput.\ Biol.\ Med.} (2024).

\bibitem{shen2025ddiben}
Y.\ Shen, X.\ Zeng, T.\ Ma, et al.,
Benchmarking drug--drug interaction prediction methods: a perspective of distribution changes,
\emph{Bioinformatics} 41 (11) (2025) btaf569.

\bibitem{openddi2025}
OpenDDI Consortium,
OpenDDI: a comprehensive benchmark for DDI prediction,
in: \emph{ICML}, 2025.

\bibitem{lukashina2021coldstart}
N.\ Lukashina, E.\ Kartysheva, O.\ Spjuth, et al.,
Cold-start problems in data-driven prediction of drug--drug interaction effects,
\emph{Pharmaceutics} 13 (9) (2021) 1400.

\bibitem{ssfddi2024}
Z.\ Zhang, Y.\ Chen, S.\ Li, et al.,
SSF-DDI: a deep learning method utilizing drug sequence and substructure features for DDI prediction,
\emph{BMC Bioinform.} 25 (2024) 38.

\bibitem{hdnddi2025}
Y.\ Wang, X.\ Li, Z.\ Chen, et al.,
HDN-DDI: a novel framework for predicting DDIs using hierarchical molecular graphs,
\emph{BMC Bioinform.} 26 (2025) 52.

\bibitem{nyamabo2021ssiddi}
A.K.\ Nyamabo, Y.\ Yu, S.\ Shi,
SSI-DDI: substructure--substructure interactions for DDI prediction,
\emph{Brief.\ Bioinform.} 22 (6) (2021) bbab133.

\bibitem{boldini2023xgboost}
D.\ Boldini, A.\ Loescher, S.\ Riniker, et al.,
Practical guidelines for the use of gradient boosting for molecular property prediction,
\emph{J.\ Cheminform.} 15 (2023) 73.

\bibitem{jiang2021gnnvsdescriptor}
D.\ Jiang, Z.\ Wu, C.Y.\ Hsieh, et al.,
Could graph neural networks learn better molecular representation for drug discovery?\ A comparison study,
\emph{J.\ Cheminform.} 13 (2021) 12.

\bibitem{shtar2019xgboost}
G.\ Shtar, L.\ Rokach, B.\ Shapira,
Detecting drug--drug interactions using artificial neural networks and classic graph similarity measures,
\emph{PLoS ONE} 14 (8) (2019) e0219796.

\bibitem{tran2021xgboost}
T.N.\ Tran, M.\ Le, T.\ Pham, et al.,
ML-based prediction of DDI for histamine antagonist using hybrid chemical features,
\emph{Cells} 10 (11) (2021) 3092.

\bibitem{sgtddi2025}
H.\ Zhu, Y.\ Liu, X.\ Wang, et al.,
SGT-DDI: spatial geometry and graph transformer for DDI prediction,
\emph{Inf.\ Fusion} (2025).

\bibitem{deng2020ddimdl}
Y.\ Deng, X.\ Xu, Y.\ Qiu, et al.,
A multimodal deep learning framework for predicting drug--drug interaction events,
\emph{Bioinformatics} 36 (15) (2020) 4316--4322.

\bibitem{mlfgnn2025}
Multi-Level Fusion GNN Authors,
Multi-level fusion graph neural network for molecule property prediction,
arXiv preprint 2507.03430, 2025.

\bibitem{zagidullin2021fingerprints}
B.\ Zagidullin, Z.\ Aldahdooh, S.\ Zheng, et al.,
Comparative analysis of molecular fingerprints in prediction of drug combination effects,
\emph{Brief.\ Bioinform.} 22 (6) (2021) bbab291.

\bibitem{wang2024knowddi}
Y.\ Wang, H.\ Min, X.\ Wu, et al.,
Accurate and interpretable DDI prediction enabled by knowledge subgraph learning,
\emph{Commun.\ Med.} 4 (2024) 48.

\bibitem{lv2023meta3d}
K.\ Lv, Y.\ Cao, Z.\ Liu, et al.,
3D graph neural network with few-shot learning for predicting DDIs in scaffold-based cold start scenario,
\emph{Neural Netw.} 165 (2023) 568--579.

\bibitem{zhang2023emergnn}
Y.\ Zhang, Q.\ Gao, K.\ Chen, et al.,
Emerging drug interaction prediction enabled by a flow-based graph neural network with biomedical network,
\emph{Nat.\ Comput.\ Sci.} 3 (2023) 1023--1033.

\bibitem{dsnddi2023}
J.\ Li, K.\ Sun, Z.\ Chen, et al.,
DSN-DDI: an accurate and generalized framework for DDI prediction by dual-view representation learning,
\emph{Brief.\ Bioinform.} 24 (1) (2023) bbac597.

\bibitem{chen2023substructure}
Y.\ Chen, T.\ Ma, X.\ Yang, et al.,
Molecular substructure-based cross-dependent network for drug--drug interaction prediction,
\emph{IEEE J.\ Biomed.\ Health Inform.} 27 (11) (2023) 5547--5558.

\bibitem{pang2023pretraining}
S.\ Pang, Y.\ Zhang, T.\ Song, et al.,
Pretraining molecular and substructural encoders for predicting drug--drug interactions in cold-start scenarios,
\emph{Brief.\ Bioinform.} 24 (6) (2023) bbad399.

\bibitem{schwarz2021attentionddi}
K.\ Schwarz, A.\ Filipiak, R.\ Loo\ss, et al.,
AttentionDDI: Siamese attention-based deep learning method for DDI predictions,
\emph{BMC Bioinform.} 22 (2021) 412.

\bibitem{cnnsiam2023}
CNN-Siam Authors,
CNN-Siam: multimodal siamese CNN-based deep learning approach for DDI prediction,
\emph{BMC Bioinform.} 24 (2023) 340.

\bibitem{faers2024}
FAERS DDI Analysis,
Drug--drug interaction adverse events from FAERS database analysis,
\emph{PMC}, 2024.

\bibitem{ren2025rareddie}
Y.\ Ren, H.\ Liu, X.\ Song, et al.,
Predicting rare drug--drug interaction events with dual-granular structure-adaptive and pair variational representation,
\emph{Nat.\ Commun.} 16 (2025) 5871.

\bibitem{clinther2022ddi}
DDI Pharmacovigilance Review,
Drug--drug interactions: a pharmacovigilance road less traveled,
\emph{Clin.\ Ther.} 44 (12) (2022) 1629--1638.

\bibitem{pharmacogenomics2023equity}
Population Pharmacogenomics Consortium,
Population pharmacogenomics for health equity,
\emph{Genes} 14 (10) (2023) 1840.

\bibitem{habre2025aiadvancing}
C.\ Habre, S.\ Fakhouri, N.\ Youssef, et al.,
Advancing drug--drug interactions research: integrating AI-powered prediction, vulnerable populations, and regulatory insights,
\emph{Front.\ Pharmacol.} 16 (2025) 1632775.

\bibitem{mhafrddi2025}
MHAFR-DDI Authors,
MHAFR-DDI: multimodal hierarchical attention fusion and relation-aware architecture for DDI event prediction,
\emph{BMC Bioinform.} (2025).

\bibitem{devilddi2024}
Z.\ Yu, H.\ Huang, S.\ Gao, et al.,
Devil in the tail: a multi-view framework for drug--drug interaction with imbalanced event distribution,
\emph{Brief.\ Bioinform.} 25 (6) (2024) bbae554.

\bibitem{karim2019kgddi}
M.R.\ Karim, M.\ Cochez, J.B.\ Jares, et al.,
Drug--drug interaction prediction based on knowledge graph embeddings and convolutional-LSTM network,
\emph{BMC Bioinform.} 20 (Suppl 10) (2019) 386.

\bibitem{daba2025}
DABA-DDI Authors,
DABA-DDI: Bayesian calibration for drug--drug interaction prediction,
\emph{PeerJ Comput.\ Sci.} (2025).

\bibitem{friesacher2025calibration}
A.H.\ Friesacher, J.\ Kirchmair, C.\ Peng, et al.,
Achieving well-informed decision-making in drug discovery: a comprehensive calibration study,
\emph{J.\ Cheminform.} 17 (2025) 22.

\end{thebibliography}

\end{document}
